\expandafter\gdef\csname SNside1\endcsname{\node[anchor=north west,inner sep=0pt,text width=48mm,font=\SideFont,align=left] at ([xshift=138mm,yshift=-516.15pt]current page.north west) {\textcolor{SNink}{\textbf{课程连接 · 分层连接操作系统}}\par\vspace{1.5mm}Cache利用短期访问局部性，虚拟存储还承担地址空间、隔离和换页。两者都能体现层次思想，但缓存对象与缺失处理不同：Cache常缓存数据块，TLB缓存映射，页表记录虚页到页框关系。};}
\expandafter\gdef\csname SNside2\endcsname{\node[anchor=north west,inner sep=0pt,text width=48mm,font=\SideFont,align=left] at ([xshift=138mm,yshift=-397.86pt]current page.north west) {\textcolor{SNink}{\textbf{最小反例 · 动态不等于访问变化}}\par\vspace{1.5mm}最小反例：一个程序只反复读同一DRAM单元，芯片仍需刷新。动态描述电荷保持机制，而非程序写入频度；随机描述按地址选取能力，也不意味着无需刷新。把这两条分类轴分开。};}
\expandafter\gdef\csname SNside3\endcsname{\node[anchor=north west,inner sep=0pt,text width=48mm,font=\SideFont,align=left] at ([xshift=138mm,yshift=-95.00pt]current page.north west) {\textcolor{SNink}{\textbf{错项排除 · 逻辑地址与物理引脚}}\par\vspace{1.5mm}4M个单元需要22位逻辑地址，但行列各11位复用时可用11根地址引脚。整机总容量不改变单片组织；若问数据与地址引脚总数，再加芯片位宽，不能直接加整机地址位数。};}
\expandafter\gdef\csname SNside4\endcsname{\node[anchor=north west,inner sep=0pt,text width=48mm,font=\SideFont,align=left] at ([xshift=138mm,yshift=-95.00pt]current page.north west) {\textcolor{SNink}{\textbf{最小反例 · 已安装容量不决定MAR}}\par\vspace{1.5mm}最小反例：机器给出64MiB字节地址空间，只安装32MiB内存。MAR仍需容纳26位地址，而非按已安装容量只取25位。芯片单元地址、可编址空间、已安装容量是三个问题。};}
\expandafter\gdef\csname SNside5\endcsname{\node[anchor=north west,inner sep=0pt,text width=48mm,font=\SideFont,align=left] at ([xshift=138mm,yshift=-95.00pt]current page.north west) {\textcolor{SNink}{\textbf{不变量 · 模块是否恢复才决定冲突}}\par\vspace{1.5mm}低位交叉中连续字轮流访问模块；同一模块相隔m个启动间隔再次被访问。无冲突要求mr不小于模块周期。只比较地址模m相同，尚不足以断言两个请求发生冲突。};}
\expandafter\gdef\csname SNside6\endcsname{\node[anchor=north west,inner sep=0pt,text width=48mm,font=\SideFont,align=left] at ([xshift=138mm,yshift=-95.00pt]current page.north west) {\textcolor{SNink}{\textbf{错项排除 · 平均旋转等半圈}}\par\vspace{1.5mm}转速r转每分意味着一圈60/r秒；均匀随机角位置的平均等待为半圈，传一扇区另按一圈中该扇区比例计算。SSD没有这段机械旋转过程，不能给它机械套用半圈等待。};}
\expandafter\gdef\csname SNside7\endcsname{\node[anchor=north west,inner sep=0pt,text width=48mm,font=\SideFont,align=left] at ([xshift=138mm,yshift=-95.00pt]current page.north west) {\textcolor{SNink}{\textbf{不变量 · 时间式与事件一一对应}}\par\vspace{1.5mm}Cache时间加缺失概率乘额外代价，前提是缺失代价不含最初查找；若已含查找，应按命中和缺失完整分支加权。两种写法可等价，但混用定义会重复计费。};}
\expandafter\gdef\csname SNside8\endcsname{\node[anchor=north west,inner sep=0pt,text width=48mm,font=\SideFont,align=left] at ([xshift=138mm,yshift=-95.00pt]current page.north west) {\textcolor{SNink}{\textbf{错项排除 · 路数不是组数}}\par\vspace{1.5mm}总M行、每组K行时有M/K组，路数为K。K=1退化为直接映射，K=M退化为全相联。用这两个端点回代，能马上排除把M/K称作路数的干扰项。};}
\expandafter\gdef\csname SNside9\endcsname{\node[anchor=north west,inner sep=0pt,text width=48mm,font=\SideFont,align=left] at ([xshift=138mm,yshift=-95.00pt]current page.north west) {\textcolor{SNink}{\textbf{最小反例 · 命中会改变LRU}}\par\vspace{1.5mm}最小反例：两路一组访问A、B、A、C。LRU在第三次命中A后淘汰B，FIFO仍按装入次序淘汰A。只记录装入顺序无法模拟LRU，命中也必须更新最近使用状态。};}
\expandafter\gdef\csname SNside10\endcsname{\node[anchor=north west,inner sep=0pt,text width=48mm,font=\SideFont,align=left] at ([xshift=138mm,yshift=-95.00pt]current page.north west) {\textcolor{SNink}{\textbf{最小反例 · 映射没缓存不等于没驻留}}\par\vspace{1.5mm}最小反例：TLB未命中，但页表存在位为1。查表后即可得到页框号，无需从磁盘调页。TLB缺失只说明映射未缓存，缺页才说明合法目标页当前不驻留。\par\smallskip\begin{tikzpicture}[x=1mm,y=1mm,>=Stealth,node distance=5mm]
\node[draw=SNink!50,fill=SNblue,rounded corners=1mm,text width=34mm,align=center,inner sep=1.2mm,font=\fontsize{7.5}{9}\selectfont] (m0) at (0,0) {TLB缺失};
\node[draw=SNink!50,fill=SNblue,rounded corners=1mm,text width=34mm,align=center,inner sep=1.2mm,font=\fontsize{7.5}{9}\selectfont] (m1) at (0,-11) {页已驻留};
\node[draw=SNink!50,fill=SNblue,rounded corners=1mm,text width=34mm,align=center,inner sep=1.2mm,font=\fontsize{7.5}{9}\selectfont] (m2) at (0,-22) {查表继续};
\draw[->,SNink] (m0) -- (m1);
\draw[->,SNink] (m1) -- (m2);
\end{tikzpicture}
\par{\fontsize{7.5}{9}\selectfont TLB未命中且存在位为1，不发生缺页。}};}
\expandafter\gdef\csname SNside11\endcsname{\node[anchor=north west,inner sep=0pt,text width=48mm,font=\SideFont,align=left] at ([xshift=138mm,yshift=-232.97pt]current page.north west) {\textcolor{SNink}{\textbf{错项排除 · 数据容量之外还有标记}}\par\vspace{1.5mm}若题目求Cache总位数，数据区之外通常还需标记、有效位，以及回写所需脏位；有相联替换时还可能需替换状态。直接映射不需要LRU状态，不能不分组织一律加同样的附加位。};}
\expandafter\gdef\csname SNside12\endcsname{\node[anchor=north west,inner sep=0pt,text width=48mm,font=\SideFont,align=left] at ([xshift=138mm,yshift=-95.00pt]current page.north west) {\textcolor{SNink}{\textbf{不变量 · 段界限先约束偏移}}\par\vspace{1.5mm}段式地址包含段号与段内偏移，应先核段表和偏移是否在界限内，再形成物理地址或进入分页。扩大基址不能修复越界偏移；非法地址与合法页未驻留也不能混成同一种缺页。};}
\expandafter\gdef\csname SNside13\endcsname{\node[anchor=north west,inner sep=0pt,text width=48mm,font=\SideFont,align=left] at ([xshift=138mm,yshift=-95.00pt]current page.north west) {\textcolor{SNink}{\textbf{错项排除 · 缺页后停止本轮查数据}}\par\vspace{1.5mm}若页表判断目标合法但未驻留，本轮尚无可用页框号，不能拿旧字段继续查数据Cache。先完成缺页处理、更新映射并重执行访问。题目未给主存字节时，也不能凭路径猜出数据值。};}
\expandafter\gdef\csname SNside14\endcsname{\node[anchor=north west,inner sep=0pt,text width=48mm,font=\SideFont,align=left] at ([xshift=138mm,yshift=-95.00pt]current page.north west) {\textcolor{SNink}{\textbf{课程连接 · MMU与OS分工}}\par\vspace{1.5mm}操作系统建立映射、分配页框并处理缺页，MMU在正常访问路径上依页表与TLB实施转换和权限检查。页表由OS管理，不表示每次TLB命中都要运行一次内核软件转换。};}
\expandafter\gdef\csname SNside15\endcsname{\node[anchor=north west,inner sep=0pt,text width=48mm,font=\SideFont,align=left] at ([xshift=138mm,yshift=-95.00pt]current page.north west) {\textcolor{SNink}{\textbf{不变量 · 相联程度不直接等于延迟}}\par\vspace{1.5mm}全相联能并行比较多行，组相联只比较一个组；实际延迟还受比较器、布线和编码影响。固定实现才可比较快慢。映射限制候选位置，替换选择牺牲行，电路时序决定访问时间。};}
\expandafter\gdef\csname SNside16\endcsname{\node[anchor=north west,inner sep=0pt,text width=48mm,font=\SideFont,align=left] at ([xshift=138mm,yshift=-95.00pt]current page.north west) {\textcolor{SNink}{\textbf{不变量 · 写回与写分配独立}}\par\vspace{1.5mm}直写或回写回答命中时怎样更新下一层，分配或不分配回答缺失时是否调入块。常见搭配不是逻辑上的唯一搭配。读清题设组合，再决定缺失数、主存写次数和脏块处理。};}
\expandafter\gdef\csname SNside17\endcsname{\node[anchor=north west,inner sep=0pt,text width=48mm,font=\SideFont,align=left] at ([xshift=138mm,yshift=-140.40pt]current page.north west) {\textcolor{SNink}{\textbf{不变量 · 相联程度不直接等于延迟}}\par\vspace{1.5mm}全相联能并行比较多行，组相联只比较一个组；实际延迟还受比较器、布线和编码影响。固定实现才可比较快慢。映射限制候选位置，替换选择牺牲行，电路时序决定访问时间。};}
\expandafter\gdef\csname SNside18\endcsname{\node[anchor=north west,inner sep=0pt,text width=48mm,font=\SideFont,align=left] at ([xshift=138mm,yshift=-95.00pt]current page.north west) {\textcolor{SNink}{\textbf{错项排除 · 缺页后停止本轮查数据}}\par\vspace{1.5mm}若页表判断目标合法但未驻留，本轮尚无可用页框号，不能拿旧字段继续查数据Cache。先完成缺页处理、更新映射并重执行访问。题目未给主存字节时，也不能凭路径猜出数据值。};}
\expandafter\gdef\csname SNside19\endcsname{\node[anchor=north west,inner sep=0pt,text width=48mm,font=\SideFont,align=left] at ([xshift=138mm,yshift=-95.00pt]current page.north west) {\textcolor{SNink}{\textbf{不变量 · 相联程度不直接等于延迟}}\par\vspace{1.5mm}全相联能并行比较多行，组相联只比较一个组；实际延迟还受比较器、布线和编码影响。固定实现才可比较快慢。映射限制候选位置，替换选择牺牲行，电路时序决定访问时间。};}
\expandafter\gdef\csname SNside20\endcsname{\node[anchor=north west,inner sep=0pt,text width=48mm,font=\SideFont,align=left] at ([xshift=138mm,yshift=-95.00pt]current page.north west) {\textcolor{SNink}{\textbf{不变量 · 段界限先约束偏移}}\par\vspace{1.5mm}段式地址包含段号与段内偏移，应先核段表和偏移是否在界限内，再形成物理地址或进入分页。扩大基址不能修复越界偏移；非法地址与合法页未驻留也不能混成同一种缺页。};}
\AddToShipoutPictureFG{\begin{tikzpicture}[remember picture,overlay]\ifcsname SNside\thepage\endcsname\csname SNside\thepage\endcsname\fi\end{tikzpicture}}